\documentclass[10pt,twocolumn]{article}
\usepackage[T1]{fontenc}
\usepackage[utf8]{inputenc}
\usepackage{lmodern}
\usepackage[margin=1.8cm,top=2.4cm,bottom=2cm,columnsep=0.6cm]{geometry}
\usepackage{fancyhdr}
\usepackage{titlesec}
\usepackage{booktabs}
\usepackage{amsmath,amssymb,amsfonts}
\usepackage{microtype}
\usepackage{xcolor}
\usepackage{mdframed}
\usepackage{parskip}
\usepackage{enumitem}
\usepackage{hyperref}

\definecolor{rulered}{RGB}{180,30,30}
\definecolor{boxgray}{RGB}{245,245,245}
\definecolor{keywordgray}{RGB}{100,100,100}

\pagestyle{fancy}
\fancyhf{}
\fancyhead[L]{\textsc{\small Claude Mag}}
\fancyhead[C]{\small\textit{Machine Learning Research Digest}}
\fancyhead[R]{\small 2026-09-16}
\fancyfoot[C]{\small\thepage}
\renewcommand{\headrulewidth}{0.8pt}

\titleformat{\section}{\normalsize\bfseries\color{rulered}}{}{0pt}{}%
  [\vspace{-4pt}\color{rulered}\rule{\columnwidth}{0.4pt}]
\titlespacing{\section}{0pt}{8pt}{4pt}

\hypersetup{colorlinks=true,urlcolor=rulered,linkcolor=black}

\begin{document}
\setlength{\parindent}{0pt}
\setlength{\parskip}{4pt}

{\small\color{keywordgray}\textbf{Keywords:}
  sparse attention \textbullet{}
  KV cache compression \textbullet{}
  end-to-end training \textbullet{}
  long-context LLM \textbullet{}
  differentiable gating}\par\vspace{4pt}

{\LARGE\bfseries\raggedright
  SAS: Simple Attention Sparsification
  via End-to-End Optimization}\par\vspace{4pt}

{\large\itshape\raggedright
  Hard Top-$K$ Blocks Gradient Flow in Trainable Sparse Attention ---
  Continuous Log-Space Gates Restore Differentiability and
  Outperform Baselines Under Tight KV Budgets}\par\vspace{6pt}

{\small\textbf{By} Zhiwei Li, Lei Zhu, Hao Gu, Xiang Hu, Yan Wang,
  Haitao Mi, Sirui Han, Leo Liang, Zhijiang Guo
  (Tencent HY LLM Frontier; HKUST)}\par\vspace{2pt}

{\small\color{keywordgray}arXiv:2609.13141}
\par\vspace{2pt}\noindent{\color{rulered}\rule{\columnwidth}{0.6pt}}\vspace{6pt}

Efficient long-context inference demands reducing attention from $O(L^2)$
to sub-linear cost by retaining only a sparse subset of KV entries per
query. Trainable selectors promise to learn which tokens matter for
downstream tasks --- but a fundamental problem has gone unaddressed:
the hard Top-$K$ gate that implements sparsification is
non-differentiable, severing gradient flow from the language modeling
loss back to the selector. SAS (Simple Attention Sparsification) fixes
this with continuous log-space gates injected directly into the
attention softmax, enabling full end-to-end training with no
architectural changes.

\begin{mdframed}[backgroundcolor=boxgray,linecolor=rulered,linewidth=1pt,
    innerleftmargin=6pt,innerrightmargin=6pt,
    innertopmargin=6pt,innerbottommargin=6pt]
\textbf{\small AT A GLANCE}\par\vspace{4pt}
\begin{description}[leftmargin=0pt,labelsep=4pt,itemsep=2pt,parsep=0pt,
    font=\small\bfseries]
  \item[Problem:] Hard Top-$K$ in trainable sparse attention blocks
    gradient flow; selectors train on surrogate losses misaligned
    with language modeling.
  \item[Why it matters:] Long-context LLMs need 5--10\% KV budgets
    to fit in memory; heuristic selectors lose 5--7 points on
    reasoning at these budgets.
  \item[Method:] Inject gate $g_i = \sigma(s_i)$ as $\log g_i$ into
    attention logits before softmax; normalize against mandatory
    current block.
  \item[Result:] Closes most of the gap to full attention at 10\%
    KV budget on reasoning, long-context, and agentic benchmarks.
  \item[Evidence:] LLaMA-3-8B and 70B, Mistral-7B; MATH-500,
    LongBench, WebArena, SWE-bench-lite.
\end{description}
\end{mdframed}

\section{Why Existing Approaches Fall Short}

Sparse attention methods select a subset $\mathcal{S} \subset
\{1,\ldots,L\}$ of KV entries per query using a lightweight selector
that outputs relevance scores $s_1,\ldots,s_L$, then applies:
\[
  \mathcal{S} = \mathrm{Top}\text{-}K(s_1,\ldots,s_L)
\]
The Top-$K$ operation has zero gradient almost everywhere:
$\partial\mathcal{S}/\partial s_i = 0$. Selectors are therefore trained
with surrogate losses such as cosine similarity between selector output
and the query --- a proxy for what matters in generation. This mismatch
is the root cause of quality degradation under tight budgets.

\section{Core Mechanism: Log-Space Gate Injection}

SAS replaces the hard gate with a continuous sigmoid gate:
\[
  g_i = \sigma(s_i) = \frac{1}{1 + e^{-s_i}} \in (0,1)
\]
The gate is injected into the attention logit before softmax:
\[
  \tilde{a}_{ij} = \mathrm{softmax}\!\left(
    \frac{\mathbf{q}_i\cdot\mathbf{k}_j}{\sqrt{d}} + \log g_j
  \right)_j
\]
Adding $\log g_j$ is equivalent to multiplying the pre-softmax
probability by $g_j$:
\[
  \tilde{a}_{ij} \propto \exp\!\left(\frac{\mathbf{q}_i\cdot\mathbf{k}_j}
  {\sqrt{d}}\right)\cdot g_j
\]
When $g_j\to 0$ a token is suppressed; when $g_j\to 1$ it is fully
retained. Because softmax is differentiable, the gradient of the
language modeling loss flows back through $\tilde{a}_{ij}$, through
the gate $g_j$, and to the selector parameters:
\[
  \frac{\partial\mathcal{L}_{\mathrm{LM}}}{\partial s_j}
  = \frac{\partial\mathcal{L}_{\mathrm{LM}}}{\partial\tilde{a}_{ij}}
    \cdot\frac{\partial\tilde{a}_{ij}}{\partial g_j}
    \cdot\sigma'(s_j)
\]

\section{Technical Formulation}

\textbf{Gate-injected attention:}
\[
  \tilde{\mathbf{a}}_i = \mathrm{softmax}\!\left(
    \frac{\mathbf{q}_i K^\top}{\sqrt{d}} + \log\mathbf{g}
  \right),\quad \mathbf{g} = \sigma(\mathbf{s})
\]

\textbf{Selector:} A 2-layer MLP or single attention head:
\[
  s_j = f_{\theta_s}(\mathbf{k}_j)
\]

\textbf{Training:} Standard causal language modeling with gradients
flowing through attention and the gate to $\theta_s$:
\[
  \mathcal{L} = -\sum_t\log p_\theta(x_t\mid x_{<t})
\]

\textbf{Inference:} Select top-$K$ entries by $s_j$ (hard), attend
only to those. The selector has been trained with soft gradients.

\section{Normalized Softmax Gates}

The current block $\mathcal{C}$ (the most recent $w$ tokens) is always
retained in full. SAS normalizes historical gate scores against the
attention mass consumed by $\mathcal{C}$:
\[
  \hat{g}_j = g_j\cdot(1 - A_{\mathcal{C}}),
  \quad A_{\mathcal{C}} = \sum_{j\in\mathcal{C}}\tilde{a}_{ij},
  \quad j\notin\mathcal{C}
\]
This calibrates the historical context budget to the remaining
attention mass after the current block is satisfied, keeping
sparsification levels consistent across sequence positions.

\section{Learning or Inference Procedure}

SAS is applied as a post-training fine-tuning step on a frozen or
lightly-tuned backbone. The selector is initialized randomly or from
heuristic scores and fine-tuned for as few as 500 steps on standard
language modeling corpora. The MLP selector operates on cached keys
and adds $O(L)$ overhead --- negligible against $O(L^2)$ attention.

\section{Experimental Results}

Evaluated on LLaMA-3-8B, LLaMA-3-70B (long-context variants) and
Mistral-7B-v0.3 at 10\% KV budget:

\begin{center}
\begin{tabular}{lccc}
\toprule
Method & MATH-500 & LongBench & WebArena \\
\midrule
Full Attention  & 68.4 & 54.2 & 17.3 \\
Quest           & 61.2 & 49.1 & 14.6 \\
SnapKV          & 59.7 & 47.8 & 13.9 \\
MagicPIG        & 63.1 & 50.4 & 15.1 \\
\textbf{SAS}    & \textbf{66.8} & \textbf{53.1} & \textbf{16.9} \\
\bottomrule
\end{tabular}
\end{center}

SAS closes most of the gap to full attention at 10\% KV budget, while
prior trainable methods lose 5--7 points on MATH-500.

\section{Ablations and Interpretation}

\textbf{Without log-gate (hard Top-$K$, end-to-end attempt):}
Training fails to converge; gate gradient is zero; performance equals
Quest level.

\textbf{Without normalized gate:} Degradation of 3.2 points on
LongBench; model over-attends to historical tokens and ignores
the current block.

\textbf{Selector depth:} 2-layer MLP is sufficient; deeper selectors
do not improve quality.

\textbf{Budget sensitivity:} SAS advantage is largest at tight budgets
(5--10\% KV); at 50\% all methods approach full attention.

\section*{Reference}

Zhiwei Li et al.\ \textbf{SAS: Simple Attention Sparsification via
End-to-End Optimization of Context Ranking.}
arXiv:2609.13141, September 2026.
\url{https://arxiv.org/abs/2609.13141}

\end{document}
